% 第三部分 A：讨论（由 main_zh.tex 引入）。

\section{讨论}\label{sec:discussion}

\paragraph{没有看到可见测试被利用的迹象。} 本文关心的问题是：用策略能看到的测试来奖励它，它可能学会满足这些测试，而不是解决问题。在我们的设置中，100 步之内没有看到这种迹象。训练提高了三个集合上解出题目的比例，包括策略从未训练过的两个集合。通过可见测试的程序中错误者的比例，在 KodCode val 和 HumanEval+ 上降低了约一半，在 MBPP+ 上几乎不变。抄写可见测试中数值的程序，仍和基座策略中一样少见。通过可见测试的错误程序，在训练前后都几乎从不是对测试的抄写：它们处理了示例所展示的情况，却把其他情况做错了。可见测试奖励会奖励这些程序，但基座策略本来就会写出它们，而在三个集合中的两个上，训练让它们在通过者中所占的比例变小了。

\paragraph{评判模型能做什么，实际做了什么。} 离线时，两个评判模型给一道题中通过可见测试的尝试排序时，都能比随机更好地分出也通过留出测试的那些（\S\ref{sec:res-study}），而且都能识破硬编码查表程序。在训练中，评判因子带来的变化没有超出种子之间的差异：在 KodCode val 和 HumanEval+ 上没有差别，在 MBPP+ 上一个种子有提升，另一个没有。剩下需要压制的不是对测试的抄写，而是在示例未展示的情况上出错的解答。在策略确实学会特判的情形下——测试更弱、题目更难或训练更长——评判因子是否有用，这些运行无法回答。

\paragraph{更好的测试去掉的是少量错误的通过，而不是未解出的题目。} 预言机奖励只奖励通过全部测试的程序。它最终在三个集合上都让通过可见测试的错误程序比 \texttt{vis} 更少（只有在 MBPP+ 上区间不含零），而解出题目的比例相同。
在 100 步之内，知道哪些可见测试通过是错误的，并没有教会策略解出更多题目；它教会的是在解不出时更少地通过可见测试。一个对正确性完美的评判模型作为因子，给出的恰好就是这个奖励。

\paragraph{局限。} 只有一种策略规模（4B）、一种任务格式（单个函数，最多三次尝试，以第一个失败的测试作为反馈）、100 步、一个取自第二部分的学习率，以及每种奖励两个种子；在 MBPP+ 上两个种子的标准差高达 \NUM{cgallsd_viskev_mbpp} 个百分点。区间对题目重采样，不对种子重采样。留出测试是 KodCode 和 EvalPlus 的测试，一个程序通不过它们，是按它们的规格判定为错，而这种规格可能比题目描述更严格。特判检查只能发现含有可见测试中数值的程序；不含此类数值的特判（比如按输入长度分支）检测不到。训练期间其他实验共用 GPU。
